\section{引言与问题建模}

本次客座讲座由哥伦比亚大学助理教授 Yunzhu Li 主讲。他领导着 Robotic Perception, Interaction and Learning（RPIL）实验室，研究方向聚焦在机器人学习——让机器人更好地感知和与物理世界交互。本讲系统介绍了机器人学习的核心范式：从问题建模、感知、强化学习、模型学习、模仿学习，到最新的机器人基础模型。

\begin{figure}[H]
\centering
\includegraphics[width=0.95\textwidth]{slides-images/slide-001.jpg}
\caption{讲座概览：Robot Learning 涉及问题建模、感知、强化学习、模型学习、模仿学习与基础模型等多个方面\protect\footnotemark}
\end{figure}
\footnotetext{Slides 第2页。}

\subsection{从监督学习到机器人学习}

在 CS231N 课程中，同学们已经学习了监督学习（给定输入 $X$ 和标签 $Y$，学习从 $X$ 到 $Y$ 的映射）和自监督学习（仅使用无标注数据提取隐含结构）。\textbf{机器人学习}与之有根本性的区别：

\begin{importantbox}{机器人学习的核心特征}
机器人学习是一个\textbf{序列决策问题}。机器人不仅要理解世界，还要\textbf{与世界交互}——它采取的动作会改变环境状态，环境会给出新的观测和奖励。目标是找到一系列动作，使得累积奖励最大化（或累积代价最小化）。
\end{importantbox}

\begin{figure}[H]
\centering
\includegraphics[width=0.95\textwidth]{slides-images/slide-003.jpg}
\caption{机器人学习的通用框架：Agent 接收任务目标和环境状态，输出动作，环境返回新状态和奖励\protect\footnotemark}
\end{figure}
\footnotetext{Slides 第4页。}

这一框架由四个核心要素组成：

\begin{itemize}
    \item \textbf{目标（Goal）}：任务的定义，如自然语言指令或目标函数
    \item \textbf{状态（State）}：环境的当前描述，如关节角度、图像观测等
    \item \textbf{动作（Action）}：Agent 对环境施加的操控，如力矩、位移等
    \item \textbf{奖励（Reward）}：环境对动作效果的反馈信号
\end{itemize}

\subsection{问题实例化}

\begin{figure}[H]
\centering
\includegraphics[width=0.95\textwidth]{slides-images/slide-004.jpg}
\caption{多种机器人学习问题的实例：CartPole 平衡、机器人运动、Atari 游戏、围棋等\protect\footnotemark}
\end{figure}
\footnotetext{Slides 第5页。}

同样的框架可以描述极为不同的问题：

\begin{table}[H]
\centering
\begin{tabular}{lllll}
\toprule
\textbf{场景} & \textbf{目标} & \textbf{状态} & \textbf{动作} & \textbf{奖励} \\
\midrule
CartPole & 保持杆平衡 & 角度/角速度/位置 & 水平力 & 每步杆立直=1 \\
机器人运动 & 向前行走 & 关节角度/速度 & 关节力矩 & 每步前进=1 \\
Atari 游戏 & 最高分 & 游戏画面像素 & 上下左右 & 分数增减 \\
围棋 & 赢棋 & 棋盘状态 & 落子位置 & 赢=1,输=0 \\
叠衣服 & 整齐叠放 & 多视角 RGB-D & 夹爪运动 & 叠好=1 \\
\bottomrule
\end{tabular}
\caption{不同任务场景下四要素的具体化}
\end{table}

\begin{knowledgebox}{机器人学习 vs 计算机视觉}
计算机视觉主要关注从高维输入中学习表征（如分类、检测）。机器人学习则是一个\textbf{约束优化问题}——物理世界是约束，目标函数定义在任务目标上，决策变量是动作序列。这是两者的根本区别。
\end{knowledgebox}

\begin{warningbox}{奖励设计的微妙之处}
即使是``叠衣服''这样看似简单的任务，奖励也可以有多种定义：面积最小？最平整？特定折法？不同用户有不同偏好。奖励设计（reward shaping）本身是一个充满微妙权衡的研究方向。
\end{warningbox}

\subsection{本章小结}

机器人学习的核心是序列决策问题，通过目标、状态、动作、奖励四要素建模。与标准的监督学习不同，机器人的动作会影响后续状态，形成闭环交互。

%% ========== Part 2: 机器人感知 ==========

